\SourceLeaf{253}%
\noindent die Gleichräumigkeit symmetrischer Körper zu beweisen, wogegen er
behauptete, dass ein ähnlicher Beweis sich schon in \textsc{Legendres}
Geometrie finde. Ich habe diess noch nicht verificiren können.

\Needspace{7\baselineskip}
\begin{center}
[\textsc{Gerling} \emph{an} \textsc{Gauss}. \emph{Marburg, 21. Juli 1844.}]
\end{center}

\{\mbox{.\ .\ .\ .\ .\ .\ .} Wegen der Bemerkung des Hrn. \textsc{Eisenstein} über die
Volumina symmetrischer Körper habe ich im \textsc{Legendre} nachgesehen.
Ich selbst habe noch die Ausgabe von 1806, hatte aber Gelegenheit,
\textsc{Crelles} Übersetzung (1844) der 12ten Auflage jetzt einzusehen; wo
ich allerdings gegen 1806 in diesen Materien einiges abgeändert finde. ---
\textsc{Legendre} hat den Satz quaestionis aber nicht \textls{unabhängig} von der
Exhaustionsmethode bewiesen. Er geht nemlich von dem Satz aus: Pyramiden
von gleicher Grundfläche und gleicher Höhe sind gleich. --- Diesen beweist
er nach der Exhaustionsmethode (im Wesentlichen jetzt ganz so wie mein
\textsc{Lorenz} § 330, statt dass er sich früher näher an Euklid gehalten).
Nun zerlegt er die symmetrischen Polyeder in symmetrische dreiseitige
Pyramiden, und schliesst auf ihre Gleichheit vermöge jenes Satzes. Der
Unterschied besteht also darin, dass \textsc{Legendre} eine weitere
Zerlegung in je 12 Pyramiden, die paarweise congruent sind, nicht ausführt. ---

Wie ich den Satz \textsc{Lorenz} § 330 gelernt habe, weiss ich selbst nicht
mehr. Wenn ich ihn aber auch selbstständig erfunden habe; so haben mir
jedenfalls Bemerkungen von \textsc{Pfaff} die Erfindung sehr nahe gelegt.
Ich erinnere mich namentlich, dass er mir (1809) mit Bedauern, dass man
hier nicht tiefer eindringen könne, von dem »Segnerschen Axiom« sprach.\}
